\clearpage
\section*{September 9, 2026: Reproducible parent decisions}
\textit{Agent-drafted implementation record. Jacob requested the committed manual source task and incorporation of the reviewed companions.}

The new \path{tasks/parent_opportunities_manual} task holds four small, versioned source tables. Its two-line Makefile requires the files without generating them. The parent producer reads symlinked inputs, preserves the prior shared-site and filing-role decisions, and applies seven additional accepted links and three rejected links. It checks every decision against final parent membership, including indirect connections. Nearby ownership matches do not become automatic acceptances.

Third Avenue now contains three additive buildings at 99 units each, totaling 297 under the documented proposed-design measure. The three source rows identify the October 2025 plans. Saved DOB I1 values remain 98, 97 and 99. The documented schedule is checked against the selected units; it is not relabeled as a first-filed or approved schedule.

The historical filing universe gains exactly one row: 3 West's 150 units. All previously included filings retain their unit values. Missing lagged land information remains missing, and the 822-unit Wharf parent is retained but ineligible for specifications requiring complete site characteristics. The land-feature producer now allows this documented missing match and reports it explicitly. Saved Attorney General searches are attached through the filing ID and exact queried address, rather than a retired parent ID; genuinely absent searches remain incomplete.

The renewed East 232nd review supports separate parents. A contemporaneous lot statement covers lot 1 alone and names Kol Zefi, whereas 769's filing names Young Bronx/Ben Baum and a different architect. Boone has stronger direct ownership evidence: the 2024 DEC application names Moses Freund and Chaim Wiesenfeld as Vaja's members, while the neighboring contract names AJ Boone/Apex and covers lots 3 and 5, excluding 1660's lot 1. These findings support separate development decisions without purporting to disprove every possible ultimate affiliation.

Wilson/Boston remains a provisional separation. The 2023 deed conveys current lot 6 to Boston Wilson Partners. The August 2026 Boston diagram divides current lot 1 into tentative lots 1 and 2 and excludes lot 6. The filings name different developers and architects. Neither document resolves common control or a development agreement at the March--April filing dates; the August record also postdates the saved extract. The manual rejection states this limitation explicitly.

Estimation remains on hold. Membership decisions do not settle unit timing, legal wage assessment, or all development phases. The audit still identifies six withdrawn applications counted with replacements, totaling 368 units, as an unresolved correction. Correctly retaining saved searches also brings the existing 253-unit Nostrand Land parent into the current audit; its three filings share a sponsor and architect, but independent boundary corroboration is pending. The original 76-parent reviews remain preserved and are mapped to current membership through filing IDs.

\textit{Reproduction:} Make in the parent, site-characteristics and audit tasks; data-panel targets in \path{tasks/analyze_485x_scale_shape_splitting/code}. The detailed review, including source URLs and page references, is \path{tasks/audits/audit_estimation_parent_links/report/ten_case_deep_review.md}. Earlier logbook entries record the pre-correction state and are superseded by this implementation entry where membership has changed.
